\documentclass[11pt]{article}
\usepackage{amsmath,amssymb,amsthm}
\usepackage{enumitem}

\newtheorem{lemma}{Lemma}
\newtheorem{theorem}{Theorem}
\newtheorem*{answer}{Answer}

\begin{document}

\title{On the derived unit \(F\to \mathrm R\pi_*\pi^{-1}F\)}
\author{}
\date{}
\maketitle

\begin{answer}
No. The stated quasi-isomorphism fails in general. We give a counterexample.
\end{answer}

\section{The counterexample}

Let \(P=\mathbb N=\{0,1,2,\dots\}\) with addition. It is an integral, saturated, sharp, commutative unital monoid.

For \(x=0\), we have \(P[-0]^\sharp=P\). For every \(x>0\), the initial homomorphism
\[
P\longrightarrow P[-x]^\sharp
\]
sending \(x\) to \(0\) is the trivial monoid \(0\). Indeed, if \(x=n\cdot 1\) with \(n>0\), then \(x\mapsto 0\) forces \(n(1)\mapsto 0\). In an integral sharp monoid, \(na=0\) implies \(a=0\), so \(1\mapsto 0\), and hence the whole monoid maps trivially. Thus
\[
B_P
\]
has exactly two objects up to isomorphism:
\[
p=(P\xrightarrow{\mathrm{id}}P)
\qquad\text{and}\qquad
z=(P\xrightarrow{0}0).
\]
There is a unique morphism
\[
p\longrightarrow z
\]
in \(B_P\), namely the quotient map \(P\to 0\). There is no morphism \(z\to p\) in \(B_P\), since that would require a homomorphism \(0\to P\) whose composite with \(P\to 0\) is \(\mathrm{id}_P\), impossible for \(P=\mathbb N\).

Therefore in
\[
C:=B_P^{\mathrm{op}},
\]
there is a unique non-identity morphism
\[
z\longrightarrow p.
\]

Let \(D:=A_P^{\mathrm{op}}\), and let
\[
i:C\hookrightarrow D
\]
be the inclusion. Define an abelian presheaf \(F\) on \(C\) by
\[
F(p)=\mathbb Z,\qquad F(z)=0,
\]
and let the unique morphism \(z\to p\) in \(C\) induce the zero map
\[
F(p)\longrightarrow F(z),
\]
i.e. the map \(\mathbb Z\to 0\). This is a well-defined presheaf, since there are no further non-identity morphisms to impose compatibility.

Let
\[
L:=i_!F
\]
be the left Kan extension of \(F\) along \(i\). We compute \(L(p)\), where \(p\) is regarded as an object of \(D\).

By the standard colimit formula for left Kan extension,
\[
L(d)=\operatorname*{colim}_{i(c)\to d} F(c).
\]
For \(d=p\), the comma category \((i\downarrow p)\) has two objects:
\[
(p,\mathrm{id}_p)
\qquad\text{and}\qquad
(z,u),
\]
where \(u:z\to p\) is the unique morphism in \(D\) induced by the quotient \(p\to z\) in \(B_P\). There is precisely one morphism
\[
(z,u)\longrightarrow (p,\mathrm{id}_p)
\]
in the comma category, induced by the unique morphism \(z\to p\) in \(C\). Hence the corresponding diagram of abelian groups is
\[
F(p)\xrightarrow{F(u)} F(z),
\]
that is,
\[
\mathbb Z\longrightarrow 0.
\]
The colimit of this diagram is \(0\). Therefore
\[
L(p)=0.
\]

Now \(\pi^{-1}F\) is defined as the sheafification \(L^+\) of \(L\) with respect to the birhizomic topology on \(D\). The unit of the adjunction
\[
F\longrightarrow \pi_*\pi^{-1}F
\]
evaluated at \(p\in C\) is the composite
\[
F(p)\longrightarrow L(p)\longrightarrow L^+(p).
\]
Since \(L(p)=0\), this composite is the zero map
\[
\mathbb Z\longrightarrow L^+(p).
\]
In particular, it is not injective.

Thus the natural morphism
\[
F\longrightarrow \pi_*\pi^{-1}F
\]
is not an isomorphism in degree \(0\). Since \(R^0\pi_*=\pi_*\), the derived morphism
\[
F\longrightarrow \mathrm R\pi_*\pi^{-1}F
\]
cannot be a quasi-isomorphism. Consequently, the proposed statement is false.

\end{document}